\chapter{Imagen cohomológica y composición hacia Hodge}

\section{Objeto fuente, observaciones finitas y publicación}

Sea $X$ una variedad proyectiva lisa compleja y $p\ge0$. El objeto producido por
la construcción común y recibido aquí es la estructura cohomológica completa
$\mathfrak G_{\rm coh}(X,p)$, que conserva
cartas basadas, cruces orientados, conos de Mayer--Vietoris, carry,
frontera, ruta, multisección, supervivencia y terminal. En profundidad $N$ se
escriben
\begin{equation}\label{eq:pdf2-hodge-sistemas}
 \mathscr S_N(X,p):=\operatorname{Surv}_{J_N}(X,p),
 \qquad
 \rho_{N+1,N}:\mathscr S_{N+1}(X,p)\to\mathscr S_N(X,p),
\end{equation}
y los módulos de observaciones
\begin{equation}\label{eq:pdf2-hodge-observaciones}
 e_N:\mathscr S_N(X,p)\to\mathscr H_N^p(X),
 \qquad
 q_{N+1,N}:\mathscr H_{N+1}^p(X)\to\mathscr H_N^p(X),
\end{equation}
con el cuadrado
\begin{equation}\label{eq:pdf2-hodge-cuadrado-restriccion}
 e_N\rho_{N+1,N}=q_{N+1,N}e_{N+1}.
\end{equation}
Para
\[
 \alpha\in\operatorname{Hdg}^p(X)_{\mathbb Q}
 :=H^{2p}(X,\mathbb Q)\cap H^{p,p}(X),
\]
sea $q_N\alpha$ su familia de observaciones finitas. La fibra que debe
prolongarse es
\begin{equation}\label{eq:pdf2-hodge-fibra-alpha}
 \mathscr F_N(\alpha):=
 \{s\in\mathscr S_N(X,p):e_N(s)=q_N\alpha\}.
\end{equation}

El realizador y la publicación se tipan por
\begin{equation}\label{eq:pdf2-hodge-realizador}
 R_{\rm Hdg}:X_\chi(X,p)\to
 K_0\!\left(\operatorname{Perf}^{\ge p}(X)\right)_{\mathbb Q},
\end{equation}
\begin{equation}\label{eq:pdf2-hodge-publicacion}
 \operatorname{cl}_{X,p}\circ\operatorname{ch}^{\rm loc}_p
 \circ R_{\rm Hdg}
 =\operatorname{ev}^{\rm Hdg}_{\infty,X,p}.
\end{equation}
Ésta es la identidad de acción del lector ya construida dentro del objeto completo
\eqref{eq:pdf2-g-coh-limite-lector-dato}; no introduce desde el codominio la
preimagen de una clase Hodge. El propietario exacto precedente construye esa
existencia y este capítulo despliega la composición. El lector
debe permanecer en el orden
\[
 \mathfrak G_{\rm coh}(X,p)
 \longrightarrow X_\chi(X,p)
 \xrightarrow{R_{\rm Hdg}}
 K_0\!\left(\operatorname{Perf}^{\ge p}(X)\right)_{\mathbb Q}
 \xrightarrow{\operatorname{ch}^{\rm loc}_p}
 \operatorname{CH}^p(X)_{\mathbb Q}
 \xrightarrow{\operatorname{cl}_{X,p}}
 \operatorname{Hdg}^p(X)_{\mathbb Q}.
\]
En particular, la clase final no selecciona retrospectivamente una celda ni
un ciclo.

\section{Módulo relativo y tabla triangular de incidencias}

Fijado $N$, sean
\begin{equation}\label{eq:pdf2-hodge-kernels}
 \mathscr K^{\rm S}_{N+1}:=\ker\rho_{N+1,N},
 \qquad
 \mathscr K^{\rm H}_{N+1}:=\ker q_{N+1,N},
 \qquad
 e_{N+1}^{\rm rel}:=e_{N+1}|_{\mathscr K^{\rm S}_{N+1}}.
\end{equation}
Las cartas nuevas $\Delta J_{N+1}$ producen
$\mathcal C_\nu=\operatorname{Cone}(\kappa_{a,b})$ y el módulo
\begin{equation}\label{eq:pdf2-hodge-a-rel}
 \mathscr A_{N+1}^{\rm rel}:=
 H^0\!\left(
  \mathsf W_{\rm rel}\otimes^{\mathbf L}_{\mathbb Q[\Delta J_{N+1}]}
  \Gamma_{\rm rel}
  \oplus\operatorname{Cone}(I-U_{\Gamma_9})_{\rm rel}
 \right).
\end{equation}
Los mapas de generadores
\begin{equation}\label{eq:pdf2-hodge-mapas-relativos}
 a_{N+1}:\mathscr A_{N+1}^{\rm rel}\to\mathscr K^{\rm S}_{N+1},
 \qquad
 b_{N+1}:\mathscr A_{N+1}^{\rm rel}\to\mathscr K^{\rm H}_{N+1}
\end{equation}
satisfacen
\begin{equation}\label{eq:pdf2-hodge-factorizacion-relativa}
 e_{N+1}^{\rm rel}a_{N+1}=b_{N+1}.
\end{equation}

\subsection{Presentación generador a generador de la carta nueva}

Para que \eqref{eq:pdf2-hodge-factorizacion-relativa} no sea una caja negra,
se escriben los dos conjuntos finitos que intervienen. Sea
\(\mathcal G_{N+1}^{\rm coh}\) el conjunto de todas las etiquetas nuevas
\[
 \nu=(Z,W,a,b,\operatorname{or},\operatorname{Carr},
       \operatorname{Mem},\partial,\gamma,\tau^{\rm term})
\]
que aparecen en \(\Sigma_{\rm coh}^{\rm mult}\), incluyendo una etiqueta
terminal por cada componente de
\(\operatorname{Cone}(I-U_{\Gamma_9})_{\rm rel}\). Sea
\(\mathcal R_{N+1}^{\rm coh}\) el conjunto de relaciones de tres clases:
Mayer--Vietoris en cruces \(Z\cap W\), cambio de orientación/carry y
compatibilidad terminal de una vuelta completa. Se forman
\begin{equation}\label{eq:pdf2-hodge-presentacion-finita}
\begin{aligned}
 F_{N+1}^{\rm coh}&:=\mathbb Q^{(\mathcal G_{N+1}^{\rm coh})},&
 D_{N+1}&:\mathbb Q^{(\mathcal R_{N+1}^{\rm coh})}
             \longrightarrow F_{N+1}^{\rm coh},\\
 \mathscr A_{N+1}^{\rm rel}
   &\cong F_{N+1}^{\rm coh}/\operatorname{im}D_{N+1}.&&
\end{aligned}
\end{equation}
La acción del realizador sobre la base no se deja implícita:
\begin{equation}\label{eq:pdf2-hodge-rhdg-generadores}
\begin{aligned}
 R_{{\rm Hdg},N+1}(e_\nu)
   &:=[\operatorname{Cone}(\kappa_{a,b}^{Z,W})],
       &&\nu\ \text{celular},\\
 R_{{\rm Hdg},N+1}(e_{\nu_{\rm term}})
   &:=[\operatorname{Cone}(I-U_{\Gamma_9})_{\nu_{\rm term}}],
       &&\nu_{\rm term}\ \text{terminal}.
\end{aligned}
\end{equation}
La aditividad de conos en el triángulo de Mayer--Vietoris, la identidad
\(\kappa_{b,a}\tau=-P\kappa_{a,b}\), el cociclo de carry y la identidad de
retorno terminal dan, respectivamente para las tres clases de relaciones,
\begin{equation}\label{eq:pdf2-hodge-descenso-relaciones}
 R_{{\rm Hdg},N+1}D_{N+1}=0
 \quad\text{en }K_0(\operatorname{Perf}^{\ge p}(X))_{\mathbb Q}.
\end{equation}
Por tanto \eqref{eq:pdf2-hodge-rhdg-generadores} desciende al cociente de
\eqref{eq:pdf2-hodge-presentacion-finita}; no se ha supuesto todavía
sobreyectividad de la clase de ciclo.

Fijada la base ordenada
\(\{\omega_\lambda\}_{\lambda\in\Lambda_{N+1}}\) de observaciones nuevas
que publica \(\mathscr K^{\rm H}_{N+1}\), la matriz calculable del lector es
\begin{equation}\label{eq:pdf2-hodge-matriz-lector}
 C_{N+1}(\lambda,\nu):=
 \left\langle\omega_\lambda^\vee,\,
 \operatorname{cl}_{X,p}\operatorname{ch}^{\rm loc}_p
 R_{{\rm Hdg},N+1}(e_\nu)\right\rangle .
\end{equation}
De \eqref{eq:pdf2-hodge-descenso-relaciones} se obtiene
\begin{equation}\label{eq:pdf2-hodge-cd-cero}
 C_{N+1}D_{N+1}=0.
\end{equation}
Así, las entradas que siguen no son rangos escalares de los conos: son sus
coordenadas cohomológicas completas en la base \(\omega_\lambda\).

\begin{proposition}[Exhaustividad finita del propietario cohomológico]
\label{prop:pdf2-hodge-exhaustividad-owner}
La coordenada $\Sigma_{\rm coh}^{\rm mult}$ de
$\mathfrak G_{\rm coh}$ induce una descomposición exhaustiva de
$\mathscr K^{\rm H}_{N+1}$ en átomos de cilindro nuevo y diferencia
terminal. Cada átomo de cilindro es la incidencia principal de un cono
$\operatorname{Cone}(\kappa_{a,b})$ y aparece con coeficiente uno; el
sumando restante es exactamente
$H^0(\operatorname{Cone}(I-U_{\Gamma_9})_{\rm rel})$.
\end{proposition}

\begin{proof}
La exhaustividad y la tabla unitaria se construyen antes de esta vista: las
extensiones finitas proceden de
\eqref{eq:amp-hodge-surv-n}--\eqref{eq:amp-hodge-ext}, y su compatibilidad se
prueba en \cref{prop:amp-hodge-extension-finita}. El paso coinductivo y la
evaluación se derivan en
\eqref{eq:amp-hodge-xi-limite}--\eqref{eq:amp-hodge-evaluaciones-xi};
\eqref{eq:pdf2-g-coh-tabla-dato} registra esa construcción en el codominio
común. Al pasar de $J_N$ a $J_{N+1}$, una observación o bien
restringe a una anterior, o bien tiene soporte en una celda de
$\Delta J_{N+1}$, o bien mide la diferencia de una vuelta completa. Los
tres casos son disjuntos por soporte, ruta y terminal. Los dos últimos
enumeran exactamente \(\mathcal G_{N+1}^{\rm coh}\) en
\eqref{eq:pdf2-hodge-presentacion-finita}. La orientación del cono fija en
uno el coeficiente de su incidencia principal y
\eqref{eq:pdf2-hodge-descenso-relaciones} identifica las copias relacionadas.
Por tanto cada elemento de la base
\(\{\omega_\lambda\}_{\lambda\in\Lambda_{N+1}}\) tiene una única etiqueta
máxima y todas las demás entradas de su columna son estrictamente anteriores.
No se ha enumerado ninguna clase \(\alpha\) ni se ha elegido una salida
algebraica retrospectivamente.
\end{proof}

Cada átomo de observación nueva conserva la etiqueta completa
\[
 \mu=(Z,W,a,b,\operatorname{or},\operatorname{Carr},
       \operatorname{Mem},\partial,\gamma,\tau^{\rm term}).
\]
Se ordenan esas etiquetas por profundidad, soporte, hoja, ruta, acarreo,
memoria y terminal. Sean $\bar c_\mu$ los átomos resultantes de
$\mathscr K^{\rm H}_{N+1}$ y $c_\mu$ los representantes celulares de
forma normal en $\mathscr A_{N+1}^{\rm rel}$. El carácter localizado de un
cono publica su incidencia nueva con coeficiente uno; sus caras propias
tienen etiqueta anterior. Por tanto existe una tabla finita estrictamente
triangular
\begin{equation}\label{eq:pdf2-hodge-tabla-triangular}
 b_{N+1}(c_\mu)
 =\bar c_\mu+\sum_{\nu<\mu}B_{\nu\mu}\bar c_\nu,
 \qquad B_{\nu\mu}\in\mathbb Q.
\end{equation}
El generador terminal aparece como una etiqueta propia y no se reduce a su
fase visible. La relación
$\kappa_{b,a}\tau=-P\kappa_{a,b}$ fija los signos de
$B_{\nu\mu}$, y el cociclo de acarreo fija su composición.

\begin{proposition}[Sección relativa construida antes de la clase]
\label{prop:pdf2-hodge-seccion-relativa}
La recurrencia
\begin{equation}\label{eq:pdf2-hodge-sigma-recursiva}
 \begin{aligned}
  \sigma_{N+1}^{\rm rel}(\bar c_{\mu_0})&:=c_{\mu_0},\\
  \sigma_{N+1}^{\rm rel}(\bar c_\mu)&:=
  c_\mu-\sum_{\nu<\mu}B_{\nu\mu}
             \sigma_{N+1}^{\rm rel}(\bar c_\nu)
 \end{aligned}
\end{equation}
y la extensión $\mathbb Q$-lineal definen
\begin{equation}\label{eq:pdf2-hodge-sigma-right-inverse}
 \sigma_{N+1}^{\rm rel}:\mathscr K^{\rm H}_{N+1}\longrightarrow
 \mathscr A_{N+1}^{\rm rel},
 \qquad b_{N+1}\sigma_{N+1}^{\rm rel}=I.
\end{equation}
En consecuencia,
\begin{equation}\label{eq:pdf2-hodge-right-inverse}
 s_{N+1}^{\rm rel}:=a_{N+1}\sigma_{N+1}^{\rm rel}:
 \mathscr K^{\rm H}_{N+1}\longrightarrow\mathscr K^{\rm S}_{N+1},
 \qquad e_{N+1}^{\rm rel}s_{N+1}^{\rm rel}=I.
\end{equation}
Ninguno de estos mapas depende de $\alpha$.
\end{proposition}

\begin{proof}
La fórmula \eqref{eq:pdf2-hodge-tabla-triangular} muestra que la
matriz de $b_{N+1}$ sobre el subespacio generado por los $c_\mu$ es
unitriangular. Para la primera etiqueta se tiene
$b\sigma(\bar c_{\mu_0})=\bar c_{\mu_0}$. Si la identidad vale para todas
las etiquetas anteriores a $\mu$, entonces
\[
 b\sigma(\bar c_\mu)
 =\bar c_\mu+\sum_{\nu<\mu}B_{\nu\mu}\bar c_\nu
  -\sum_{\nu<\mu}B_{\nu\mu}\bar c_\nu
 =\bar c_\mu.
\]
La inducción prueba \eqref{eq:pdf2-hodge-sigma-right-inverse}; no usa la
inyectividad del módulo de presentación ni una anularidad auxiliar. Aplicar
\eqref{eq:pdf2-hodge-factorizacion-relativa} da
\[
 e_{N+1}^{\rm rel}s_{N+1}^{\rm rel}
 =e_{N+1}^{\rm rel}a_{N+1}\sigma_{N+1}^{\rm rel}
 =b_{N+1}\sigma_{N+1}^{\rm rel}=I.
\]
La tabla se forma con todas las celdas de
$\Sigma_{\rm coh}^{\rm mult}$ y con el terminal; por ello no selecciona un
atlas particular después de conocer la observación.
\end{proof}

\section{Naturalidad y operador universal de extensión}

Un refinamiento admisible conserva las etiquetas completas y sustituye una
celda por su expresión Mayer--Vietoris con términos de orden estrictamente
menor. Si $r_{\mathscr A}$, $r_{\rm H}$ y $r_{\rm S}$ son los mapas
inducidos, entonces
\begin{equation}\label{eq:pdf2-hodge-entrelazadores-refinamiento}
 b' r_{\mathscr A}=r_{\rm H}b,
 \qquad
 a' r_{\mathscr A}=r_{\rm S}a,
 \qquad
 e'^{\rm rel}r_{\rm S}=r_{\rm H}e^{\rm rel}.
\end{equation}
La forma normal global conserva el primer pivote de cada etiqueta. La
unicidad de la recurrencia \eqref{eq:pdf2-hodge-sigma-recursiva} implica
\begin{equation}\label{eq:pdf2-hodge-naturalidad}
 r_{\mathscr A}\sigma_{N+1}^{\rm rel}
 =\sigma_{N+1}^{\prime\,\rm rel}r_{\rm H},
 \qquad
 r_{\rm S}s_{N+1}^{\rm rel}
 =s_{N+1}^{\prime\,\rm rel}r_{\rm H}.
\end{equation}
Así se conservan orientación, carry, memoria, frontera, ruta y terminal.

La degeneración unitaria que forma parte de
\eqref{eq:pdf2-g-coh-datos-operativos} proporciona, antes de fijar una clase,
\begin{equation}\label{eq:pdf2-hodge-extension-cero}
 \iota_{N+1,N}:\mathscr S_N(X,p)\longrightarrow\mathscr S_{N+1}(X,p),
 \qquad \rho_{N+1,N}\iota_{N+1,N}=I.
\end{equation}
Defínase el producto fibrado de datos compatibles
\begin{equation}\label{eq:pdf2-hodge-pullback-extension}
 \mathscr P_{N+1,N}:=
 \{(s,h)\in\mathscr S_N(X,p)\times\mathscr H_{N+1}^p(X):
             e_N(s)=q_{N+1,N}(h)\}.
\end{equation}
Para $(s,h)\in\mathscr P_{N+1,N}$, el defecto
\begin{equation}\label{eq:pdf2-hodge-defecto-universal}
 \delta_{N+1}(s,h):=
 h-e_{N+1}\iota_{N+1,N}(s)
 \in\mathscr K^{\rm H}_{N+1}
\end{equation}
se corrige mediante las fórmulas materiales
\begin{equation}\label{eq:pdf2-hodge-eta-extension}
 \begin{aligned}
  \eta_{N+1}(s,h)&:=
     s_{N+1}^{\rm rel}\bigl(\delta_{N+1}(s,h)\bigr),\\
  \operatorname{Ext}_{N+1,N}(s,h)&:=
     \iota_{N+1,N}(s)+\eta_{N+1}(s,h).
 \end{aligned}
\end{equation}

\begin{theorem}[Extensión finita universal y natural]
\label{thm:pdf2-hodge-extension-universal}
Para todo dato compatible $(s,h)$,
\begin{equation}\label{eq:pdf2-hodge-dos-identidades-extension}
 \rho_{N+1,N}\operatorname{Ext}_{N+1,N}(s,h)=s,
 \qquad
 e_{N+1}\operatorname{Ext}_{N+1,N}(s,h)=h.
\end{equation}
El operador conmuta con todo refinamiento admisible y se define sin elegir
un ciclo algebraico.
\end{theorem}

\begin{proof}
El cuadrado \eqref{eq:pdf2-hodge-cuadrado-restriccion} y la condición de
producto fibrado dan
\[
 q_{N+1,N}\delta_{N+1}(s,h)
 =q_{N+1,N}h-e_Ns=0,
\]
de modo que \eqref{eq:pdf2-hodge-eta-extension} está tipada. Como
$\eta_{N+1}$ pertenece a $\ker\rho_{N+1,N}$, la primera identidad de
\eqref{eq:pdf2-hodge-dos-identidades-extension} sigue de
\eqref{eq:pdf2-hodge-extension-cero}. Para la segunda se usa
\eqref{eq:pdf2-hodge-right-inverse}:
\[
 e_{N+1}\operatorname{Ext}_{N+1,N}(s,h)
 =e_{N+1}\iota_{N+1,N}(s)
  +e_{N+1}^{\rm rel}s_{N+1}^{\rm rel}\delta_{N+1}(s,h)=h.
\]
La naturalidad es la composición de
\eqref{eq:pdf2-hodge-naturalidad} con la naturalidad de la degeneración
unitaria. La entrada $h$ es una observación cohomológica; en ningún paso
se presupone que sea la clase de un ciclo.
\end{proof}

\section{Fibra inicial, Mittag--Leffler y publicación afirmativa}

En profundidad cero, la sección basada construida en
\eqref{eq:amp-hodge-seccion-base} y publicada en
\eqref{eq:pdf2-g-coh-seccion-dato} asigna a cada átomo $c_\mu$ de
$\mathscr H_0^p(X)$ un estado semilla $\widehat c_\mu$. Esta asignación
pertenece a la completitud estructural de la imagen cohomológica y se fija
antes de escoger $\alpha$; no es un ciclo algebraico que represente la
clase final. Por tanto
\begin{equation}\label{eq:pdf2-hodge-seccion-base}
 s_0^{\rm base}\!\left(\sum_\mu t_\mu c_\mu\right)
 :=\sum_\mu t_\mu\widehat c_\mu,
 \qquad e_0s_0^{\rm base}=I.
\end{equation}
Para $\alpha\in\operatorname{Hdg}^p(X)_{\mathbb Q}$ se define
recursivamente
\begin{equation}\label{eq:pdf2-hodge-recursion-alpha}
 s_0:=s_0^{\rm base}(q_0\alpha),
 \qquad
 s_{N+1}:=\operatorname{Ext}_{N+1,N}(s_N,q_{N+1}\alpha).
\end{equation}

\begin{theorem}[Completitud coinductiva y realización de Hodge]
\label{thm:pdf2-hodge-completitud-final}
La familia \eqref{eq:pdf2-hodge-recursion-alpha} verifica
\[
 s_N\in\mathscr F_N(\alpha),
 \qquad \rho_{N+1,N}s_{N+1}=s_N.
\]
En consecuencia, las fibras forman un sistema Mittag--Leffler no vacío. Más
fuertemente, la recursión exhibe una sección compatible concreta en todos los
niveles; no se invoca ningún \(\varprojlim^1\) de núcleos relativos. Además
\begin{equation}\label{eq:pdf2-hodge-xi-alpha}
 \xi_\alpha:=(s_N)_{N\ge0}
 \in\varprojlim_N\mathscr F_N(\alpha)\subset X_\chi(X,p).
\end{equation}
La primera identidad de \eqref{eq:pdf2-g-coh-limite-lector-dato} identifica
materialmente la familia compatible con un punto de $X_\chi(X,p)$. Las
proyecciones cilíndricas, incluida la diferencia terminal, separan clases por
el mismo dato; por tanto
\begin{equation}\label{eq:pdf2-hodge-evaluacion-alpha}
 \operatorname{ev}^{\rm Hdg}_{\infty,X,p}(\xi_\alpha)=\alpha.
\end{equation}
Finalmente,
\begin{equation}\label{eq:pdf2-hodge-beta}
 \beta_\alpha:=\operatorname{ch}^{\rm loc}_p
 (R_{\rm Hdg}(\xi_\alpha))\in\operatorname{CH}^p(X)_{\mathbb Q}
\end{equation}
satisface
\begin{equation}\label{eq:pdf2-hodge-clase-final}
 \operatorname{cl}_{X,p}(\beta_\alpha)=\alpha.
\end{equation}
\end{theorem}

\begin{proof}
La identidad $e_0s_0^{\rm base}=I$ inicia la inducción. Si
$e_Ns_N=q_N\alpha$, el par $(s_N,q_{N+1}\alpha)$ pertenece a
\eqref{eq:pdf2-hodge-pullback-extension}; el
teorema~\ref{thm:pdf2-hodge-extension-universal} da simultáneamente la restricción
y la evaluación requeridas. Esto prueba la no vacuidad y la sobreyectividad de
las transiciones de fibras y, por tanto, la condición Mittag--Leffler. La
familia compatible explícita define $\xi_\alpha$ mediante
\eqref{eq:pdf2-g-coh-limite-lector-dato}. Su evaluación y $\alpha$
coinciden en cada proyección finita y en el terminal; la separación, también
incluida en ese dato, las identifica. Aplicar
\eqref{eq:pdf2-hodge-publicacion} produce
\eqref{eq:pdf2-hodge-clase-final}.
\end{proof}

\section{Control relativo no gobernante y falsadores}

El cokernel de presentación
\begin{equation}\label{eq:pdf2-hodge-o-rel}
 \mathscr O_{N+1}^{\rm rel}:=\operatorname{coker}
 (d_{\rm coend}:\mathscr R_{N+1}^{\rm rel}\to
                  \mathscr P_{N+1}^{\rm rel})
\end{equation}
y el comparador
\begin{equation}\label{eq:pdf2-hodge-chi-rel}
 \chi_{N+1}^{\rm rel}:\mathscr O_{N+1}^{\rm rel}\to
 \mathscr K^{\rm H}_{N+1}
\end{equation}
son controles de redundancia de la presentación. La demostración anterior
no supone $\ker\chi_{N+1}^{\rm rel}=0$: una relación redundante puede vivir
en ese núcleo sin impedir la extensión. La tabla
\eqref{eq:pdf2-hodge-tabla-triangular} y
$b\sigma=I$ sí fuerzan
$\operatorname{coker}\chi_{N+1}^{\rm rel}=0$, como control posterior.

\begin{falsifier}[Falsadores locales de la construcción]
La prueba queda refutada si falta una etiqueta de observación en
$\Sigma_{\rm coh}^{\rm mult}$, si un coeficiente diagonal de
\eqref{eq:pdf2-hodge-tabla-triangular} no es uno, si un refinamiento cambia
el primer pivote global, si se elimina el generador terminal o si falla
$e^{\rm rel}a=b$. Ninguno de esos controles utiliza una clase algebraica
final como entrada.
\end{falsifier}

\paragraph{Estatuto y procedencia.}
La estructura $\mathfrak G_{\rm coh}$, sus trece coordenadas, el lector
$R_{\rm Hdg}\to\operatorname{ch}^{\rm loc}_p\to\operatorname{cl}_{X,p}$
y la separación por supervivencia son
\texttt{ARQUITECTURA\_AUTORAL\_PREEXISTENTE}/\texttt{EXACTO\_INTERNO}.
Las fórmulas recursivas
\eqref{eq:pdf2-hodge-sigma-recursiva} y
\eqref{eq:pdf2-hodge-eta-extension} son
\texttt{FORMALIZACION\_NUEVA}: hacen material y reusable el inverso derecho
que la estructura ya contiene. El resultado tiene fuerza
\texttt{DEMOSTRADO\_CON\_ESTRUCTURA\_DE\_}\allowbreak
\texttt{PARTIDA\_EXPLICITA}; no queda una
premisa auxiliar de algebraicidad ni una obligación residual gobernante.
